\documentclass[11pt]{article}
\usepackage[margin=1in]{geometry}
\usepackage{amsmath,amssymb}
\usepackage{booktabs}
\usepackage{array}
\usepackage{enumitem}
\usepackage{hyperref}

\title{Simulation Plan: Nonlinearity Envelopes of the LuGre--Adamov Mecanum Plant}
\author{Vishvesh Koranne}
\date{\today}

\begin{document}
\maketitle

\section{Intent (from \texttt{intent.txt})}
Measure how nonlinear the full plant is (LuGre--Adamov friction included) by
varying physically meaningful dimensionless constants, and map the result onto
classes of motion to decide where a linear controller suffices and where a
nonlinear/robust controller is warranted. The dimensionless groups are defined in
\texttt{nd\_analysis.tex}; this document plans the simulations.

\medskip
\textbf{Deliverable.} For each plant/environment set, the boundary in motion space
at which the nonlinearity measures flip from ``linear'' to ``nonlinear'', and a
table translating that boundary into motion classes (cruise, aggressive
acceleration, turning, reversal, creep/docking).

\section{What is varied}

\subsection{Plant / environment set (outer loop)}
\begin{tabular}{@{}lll@{}}
\toprule
Constant & Levels & Physical lever \\
\midrule
$\mu$ & 0.3, 0.5, 0.8 & floor surface \\
$(\varepsilon_X, \varepsilon_Y)$, $\iota$ & COM offset magnitude $\times\{0, 1, 3\}$, & payload placement \\
 & angle $\phi_c$ every $45^\circ$; $\iota$ by parallel axis & \\
$\hat\chi$ ($\kappa_\chi$) & 0.021 (PU), 0.046 (rubber est.) & contact patch \\
Material set & PU vs.\ rubber: $\hat\omega_b$ via $\sigma_0$, $\zeta_b$, $\hat\chi$, $r_s$ & roller material \\
Friction case & 1, 2 ($\beta_1$, $\beta_2$ scaled by 10) & drivetrain/bearing losses \\
\bottomrule
\end{tabular}

\medskip
\textbf{COM offset and yaw inertia move together.} Moving the payload shifts the
COM and, by the parallel-axis theorem, the yaw inertia about the geometric centre
(the reference point of the plant's mass matrix):
\begin{equation}
I_s = I_G + m\,(a_X^2 + a_Y^2), \qquad
\iota = \iota_G + r_m\,\frac{\varepsilon_X^2 + \Lambda^2\varepsilon_Y^2}{(1+\Lambda)^2},
\end{equation}
with $I_G = I_s^{\mathrm{base}} - m(a_{X0}^2 + a_{Y0}^2)$ held fixed. At $3\times$ the
base offset $I_s$ rises by about $5\%$. The offset is swept as a magnitude and a
direction $\phi_c$, so that the COM can sit toward any wheel or edge.

\medskip
Held fixed (robot design): $r_m$, $\Lambda$, $\iota_G$, $\hat\jmath_w$,
$\hat\jmath_r$, $\rho$, $\rho_a$, $\kappa_\sigma$, $\hat\tau_V$, $\hat\tau_i$,
$\hat\tau_f$. The full factorial is $3\times(1 + 2\times8)\times2\times2 = 204$ sets; the
preliminary sensitivity study (P5) is expected to prune it.

\subsection{Motion set (inner loop)}
For a given plant set, a motion is fixed by $V_0$, $\Omega_0$, $T$, the
translation direction $\beta$ in the body frame, and the waveform. In
non-dimensional form the independent motion axes are
\begin{gather*}
u = \frac{V_0}{v_{\mathrm{nl}}}\ \text{(speed vs.\ actuator)},\qquad
\nu = \frac{h\Omega_0}{V_0}\ \text{(rotation share)},\\
\Pi_a = \frac{V_0}{\mu g T}\ \text{(acceleration demand)},\qquad
\beta\ \text{(direction)},
\end{gather*}
plus the waveform class. Every other motion group in \texttt{nd\_analysis.tex}
follows from these and the plant set: e.g.\
$\Pi_c = \nu\,u^2\,v_{\mathrm{nl}}^2/(\mu g h)$, $\beta_m = \hat\tau_V\,u/\mu$,
$\Pi_b$, $\hat v$, $\hat\omega_b$. The translation direction spans the full circle,
$\beta \in [0^\circ, 360^\circ)$ in $15^\circ$ steps (24 directions), and the yaw
rate is signed ($\pm\nu$). The full circle is needed because the mecanum plant is
anisotropic (Tier C found cross-coupling near 1 at the lateral cap) and an
off-centre COM breaks the plant's mirror symmetries, so no quadrant can be
inferred from another. P3 uses $45^\circ$ steps to locate sharp features; the
main sweep refines around them.

\medskip
\textbf{Waveform classes} (each probes a different nonlinearity):
\begin{enumerate}[leftmargin=*,itemsep=2pt]
  \item \emph{Cruise + small signal}: steady $(u,\nu,\beta)$ with a small multisine
    on top. Probes the local plant (Jacobian, BLA).
  \item \emph{Aggressive transient}: ramp/step to $(u,\nu,\beta)$ in time $T$.
    Probes friction saturation via $\Pi_a$.
  \item \emph{Reversal / oscillation}: velocity sign changes. Probes the
    pre-sliding stick mode, Stribeck dip, $\hat\tau_f$.
  \item \emph{Creep / docking}: $V_0$ of order $v_{\mathrm{str}}$ and holds.
    Probes $\hat v$, $\Pi_b$, the bristle linear zone.
\end{enumerate}

\subsection{Nonlinearity measures and flip thresholds}
Reused from \texttt{nonlinearity\_analysis/} (see its \texttt{REPORT.md}):

\medskip
\begin{tabular}{@{}p{0.27\linewidth}p{0.27\linewidth}p{0.38\linewidth}@{}}
\toprule
Measure & Flip criterion & Answers \\
\midrule
Tier A: Allgöwer $\nu_A$ (static map) & feasibility only & where steady motion is reachable \\
Tier C: local Jacobian spread, cross-coupling & gain spread $> 2\times$ & how much the local linear model changes \\
Tier B: BLA distortion (odd multisine) & $> -20$\,dB & open-loop, amplitude-dependent nonlinearity \\
Tier F: Gamma $\hat\gamma$ (closed loop, linear PID) & $\hat\gamma > 1$ & whether a linear controller is enough \\
\bottomrule
\end{tabular}

\medskip
Tier F is the measure closest to the intent; Tiers B and C explain \emph{why} it
flips.

\section{Preliminary simulations (to fix the sweep ranges)}

\begin{description}[leftmargin=0pt,itemsep=6pt]
\item[P0. Toolchain and reproduction.] Add the remaining packages needed by
  \texttt{run\_one.jl} and \texttt{hybrid\_ctrl} to the folder-local environment
  (Plots, ProgressMeter, FunctionWrappers, OSQP, MatrixEquations, JSON; disk
  permitting). Reproduce the stored base Tier A+C result ($\nu_A = 0.051$,
  $\approx 16$\,s) and one Tier B case, and regenerate the missing
  \texttt{.jld2} baselines needed by Tiers B and F.
  \emph{Output:} a validated Linux toolchain.

\item[P1. Timing benchmark.] Time one evaluation of each measure (Tier C point,
  one Tier B multisine run, one Tier F window) at the base plant.
  \emph{Output:} cost per point, used to size every later grid.

\item[P2. Feasible motion envelope per plant set.] Using the Tier A
  steady-state solver and open-loop voltage ramps, find for each $\mu$ and
  $\beta$: the maximum steady $u$ (actuator limit, $v_{\mathrm{nl}}$, back-EMF),
  the maximum $\nu$, and the acceleration at which the first wheel reaches the
  friction limit (largest usable $\Pi_a$).
  \emph{Output:} upper bounds $u_{\max}(\beta,\nu;\mu)$ and $\Pi_{a,\max}$.

\item[P3. One-at-a-time transition scans.] From the base plant, sweep one motion
  axis at a time on a log grid ($\Pi_a$, $u$, $\nu$, and $V_0$ down to
  $\sim v_{\mathrm{str}}$ for creep), evaluating the cheap measures (Tier C, a few
  Tier B amplitudes). Locate where each measure starts to change, including the
  friction/actuator crossover predicted at $u \approx 1 - 0.47\mu$.
  \emph{Output:} the region of each axis where the fine grid must be dense.

\item[P4. Similarity (collapse) check.] Run two dimensionally different
  configurations with identical dimensionless groups (e.g.\ scaled $V_0$, $T$
  and $\mu$ together, or a geometrically scaled robot) and compare their
  non-dimensional responses and measures.
  \emph{Output:} confirmation that the chosen groups are sufficient; if not, the
  missing group is identified before the main sweep.

\item[P5. Plant-parameter sensitivity.] At three representative motion points
  (one inside, one near, one beyond the P3 transition) vary each plant/environment
  constant alone. Drop constants that do not move the measures or the boundary
  (the existing static scan found COM and $\hat\chi$ mostly shrink feasibility).
  \emph{Output:} the pruned outer-loop set.

\item[P6. Regularisation sensitivity.] Near stick (creep and reversal classes),
  vary $\varepsilon_{\mathrm{reg}}$ ($10^{-5}$--$10^{-3}$\,m/s), because the bristle
  linear-zone threshold is of the same order.
  \emph{Output:} whether low-speed results are physical or numerical.
\end{description}

\paragraph{Known modelling caveat (not a separate study).} The roller equation
appears to use $-M_zC$ where energy consistency requires $+M_zC$. Its size is
about 5\% of the roller torque balance and the roller is slaved, so it is
recorded as a limitation; P0 can log the power-balance residual alongside the
reproduction run at no extra cost.

\section{Main sweep (after the preliminaries)}
\begin{enumerate}[leftmargin=*,itemsep=2pt]
  \item Outer loop: the pruned plant sets from P5.
  \item Inner loop: for each $(\pm\nu,\beta)$ over the full circle and waveform class, find the flip
    boundary along $\Pi_a$ (and along $u$ for the cruise class) by bisection on
    the Tier F threshold $\hat\gamma = 1$, using the ranges from P2 and P3, instead
    of a full grid.
  \item At each boundary point evaluate Tiers B and C to attribute the flip to a
    mechanism: friction saturation, actuator/back-EMF saturation, stick mode, or
    geometric coupling.
  \item Assemble boundary surfaces $\Pi_a^*(u,\nu,\beta\,|\,\text{plant set})$ and
    the motion-class table (linear sufficient / nonlinear warranted).
\end{enumerate}

\section{Decisions still open}
\begin{itemize}[itemsep=2pt]
  \item \textbf{Plant:} voltage-input \texttt{hybrid\_ctrl/plant.jl} (recommended:
    has the back-EMF limit, the corrected yaw arm, and matches the existing
    measures) vs.\ torque-input \texttt{run\_one.jl}.
  \item \textbf{Threads:} the machine has 32 cores and 125\,GB RAM; the repo's
    $\le 8$-thread rule was written for the Windows laptop.
  \item \textbf{Linear controller for Tier F:} the archived PID-CT gains were tuned
    at $\mu = 0.5$; either keep them fixed (tests robustness of one linear design)
    or retune per $\mu$ (tests whether \emph{any} linear design suffices).
  \item \textbf{Rubber parameters:} replace the Hertz/Mindlin estimates with
    measured or cited values before P5.
  \item \textbf{Result storage:} sweep outputs written under
    \texttt{NonDimensionalAnalysis/results/} (small summaries only), with raw
    trajectories kept out of git.
\end{itemize}

\end{document}
